\subsection{Analisis Delegasi Akses Antar Tenaga Kesehatan}
\label{subsec:analisis-delegasi-atenuasi-offline}

Delegasi akses antar tenaga kesehatan diperlukan berdasarkan hasil analisis pada Subbab \ref{subsubsec:analisis-skenario-layanan-klinis}. Dalam skenario DecMed sebelumnya, pasien memberikan akses kepada aktor tertentu secara langsung. Pola tersebut kurang sesuai untuk layanan klinis kolaboratif karena kebutuhan akses dapat berpindah antar aktor selama episode berlayanan berlangsung. Jika setiap perpindahan akses harus selalu diterbitkan langsung oleh pasien, maka proses pelayanan menjadi bergantung pada ketersediaan pasien.

Inti dari fitur delegasi antar tenaga kesehatan yang dianalisis adalah bahwa proses delegasi tidak melibatkan pasien dan tidak menerbitkan token baru. Setelah pasien memberikan akses awal, aktor yang sah dapat menurunkan hak akses yang sudah diterima kepada tenaga kesehatan lain dengan menambahkan batasan yang mempersempit cakupan akses.

Alternatif solusi pertama yang dipertimbangkan adalah delegasi akses melalui \textit{client} dengan pencatatan hasil delegasi pada penyimpanan \textit{on-chain}. Pada pendekatan ini, aktor yang memiliki token induk membentuk token turunan secara lokal dengan menambahkan batasan baru terhadap hak akses yang dimiliki. Token turunan tersebut kemudian dibagikan kepada penerima delegasi melalui penyimpanan \textit{on-chain}. Pendekatan ini dapat mengurangi keterlibatan pasien karena pemegang token dapat menurunkan hak akses yang sudah diterima sebelumnya.

Namun, pendekatan tersebut menjadikan \textit{client} sebagai komponen yang menentukan hasil delegasi. Kondisi ini kurang ideal karena \textit{client} dapat dimanipulasi untuk melakukan delegasi sesuka hati. Akibatnya, sistem berpotensi mencatat delegasi yang tidak konsisten dengan hak akses induk, meskipun token turunan tersebut kemudian dapat ditolak ketika digunakan untuk mengakses \textit{Server} PRE.

Alternatif berikutnya adalah melibatkan \textit{Server} PRE sebagai komponen yang memverifikasi permintaan delegasi sebelum token turunan dibentuk. Pada pendekatan ini, aktor delegator tetap menginisiasi delegasi melalui \textit{client}. \textit{Client} mengirimkan token induk yang dimiliki \textit{delegator}, lalu token dan hak akses \textit{delegator} diverifikasi oleh \textit{Server} PRE.

Jika permintaan delegasi dinyatakan sah, \textit{Server} PRE membentuk hak akses turunan dengan menambahkan batasan baru yang mempersempit hak akses. PRE kemudian mengembalikan hak akses tersebut kepada \textit{client}. Setelah menerima hasil atenuasi dari PRE, \textit{client} kemudian dapat melakukan proses pencatatan lebih lanjut pada penyimpanan \textit{on-chain}. Dengan demikian, PRE hanya menjadi verifikator dan pembentuk hak akses turunan, sedangkan pencatatan delegasi tetap dilakukan oleh \textit{client}.

Berdasarkan pertimbangan tersebut, mekanisme delegasi akses antar tenaga kesehatan yang dimediasi oleh \textit{Server} PRE dipilih. Mekanisme ini memanfaatkan konsep atenuasi sehingga hak akses turunan hanya dapat memiliki cakupan yang sama atau lebih sempit daripada hak akses induk. Setelah pasien memberikan akses awal, aktor yang sah dapat mengajukan delegasi terbatas melalui \textit{client} tanpa melibatkan pasien kembali dan tanpa menerbitkan token baru, sedangkan PRE memverifikasi dan membentuk hak akses turunan berdasarkan hak akses induk yang masih berlaku.

Pemilihan mekanisme ini memberikan batas tanggung jawab yang lebih jelas. \textit{client} berperan sebagai inisiator delegasi dan pencatat delegasi pada penyimpanan \textit{on-chain}. \textit{Server} PRE berperan sebagai pembentuk hak akses turunan. Penyimpanan \textit{on-chain} mencatat riwayat delegasi dan bukti akses. Dengan pembagian tersebut, proses pelayanan tidak perlu selalu melibatkan pasien, tetapi pembentukan hak akses turunan tetap melewati validasi PRE sebelum dicatat dan digunakan.
